\documentclass{labreport}

\usepackage{microtype}
\usepackage{enumitem}
\usepackage{pifont}
\usepackage{tabularx}
\usepackage{booktabs}

% Headheight adjustment for fancyhdr
\setlength{\headheight}{28pt}
\addtolength{\topmargin}{-13pt}

% Colors
\definecolor{codebg}{RGB}{248, 249, 251}
\definecolor{codeborder}{RGB}{218, 222, 229}
\definecolor{linkcolor}{RGB}{20, 60, 120}
\definecolor{zshgreen}{RGB}{0, 160, 80}
\definecolor{zshred}{RGB}{200, 30, 30}

% Hyperref setup with subtle link colors
\hypersetup{
    colorlinks=true,
    linkcolor=linkcolor,
    filecolor=linkcolor,
    urlcolor=linkcolor,
    citecolor=linkcolor
}

% Code and terminal listings
\lstdefinestyle{bash}{
    basicstyle=\ttfamily\footnotesize,
    backgroundcolor=\color{codebg},
    frame=single,
    framerule=0.5pt,
    rulecolor=\color{codeborder},
    breaklines=true,
    breakatwhitespace=false,
    showstringspaces=false,
    keepspaces=true,
    columns=flexible,
    numbers=none,
    xleftmargin=3pt,
    xrightmargin=3pt,
    framesep=4pt
}

\lstdefinestyle{c}{
    language=C,
    basicstyle=\ttfamily\footnotesize,
    backgroundcolor=\color{codebg},
    frame=single,
    framerule=0.5pt,
    rulecolor=\color{codeborder},
    keywordstyle=\color{blue!85!black}\bfseries,
    commentstyle=\color{gray!80!black}\itshape,
    stringstyle=\color{teal!80!black},
    numbers=left,
    numberstyle=\tiny\color{gray!80!black},
    numbersep=8pt,
    breaklines=true,
    breakatwhitespace=false,
    showstringspaces=false,
    keepspaces=true,
    columns=flexible,
    tabsize=2,
    xleftmargin=14pt,
    xrightmargin=3pt,
    framesep=4pt
}

\lstset{
    extendedchars=true,
    literate=
        {➜}{{\textcolor{zshgreen}{\ding{220}}}}{1}
        {├─}{|-\,}{2}
        {└─}{\`-\,}{2}
        {├}{|}{1}
        {─}{-}{1}
        {└}{\`}{1}
        {✗}{{\textcolor{zshred}{\ding{55}}}}{1}
}

% Section styling for clean document structure
\titleformat{\section}
  {\large\bfseries}{\thesection}{0.8em}{}[\vspace{1pt}\hrule height 0.5pt]
\titleformat{\subsection}
  {\normalsize\bfseries}{\thesubsection}{0.8em}{}
\titleformat{\subsubsection}
  {\small\bfseries}{\thesubsubsection}{0.8em}{}

\begin{document}

\labtitle
    {13}
    {Semaphores in Your Own Kernel}
    {st. gr. FAF-243}
    {Șoimu Ionuț}
    {asist. univ.\\[0.15cm] Reițman Patricia}

\tableofcontents
\newpage

% ==============================================================================
\section{Lab 13: Track B -- Semaphores in Your Own Kernel (xv6)}

\subsection{Before you start}

\subsubsection*{Purpose and Scope}
In Laboratories 11 and 12, synchronization primitives were consumed from user space through the standard POSIX C library (\texttt{pthread\_mutex\_t}, \texttt{sem\_t}, \texttt{pthread\_cond\_t}). While these libraries appear to work effortlessly, their functionality relies on kernel-space mechanics: atomic instructions, spinlocks, and sleep/wakeup scheduling queues.

The objective of Laboratory Work 13 is to transition into **kernel space** and construct a complete, native counting semaphore implementation inside the **xv6-riscv** kernel. By implementing three foundational system calls:
\begin{enumerate}[label=(\alph*), noitemsep]
    \item \texttt{sem\_init(int id, int value)} (Syscall 24): Initializes semaphore slot \texttt{id} to \texttt{value};
    \item \texttt{sem\_down(int id)} (Syscall 25): Acquires a permit, or suspends the calling process without wasting CPU cycles;
    \item \texttt{sem\_up(int id)} (Syscall 26): Restores a permit and awakens waiting processes.
\end{enumerate}

\subsubsection*{Architectural Nuances of our xv6 Tree}
In our modern course xv6-riscv codebase:
\begin{itemize}[noitemsep]
    \item Syscall 22 is allocated to \texttt{SYS\_sync}, and syscall 23 was allocated to \texttt{SYS\_trace} in Lab 10. Consequently, the semaphore system calls occupy indices \textbf{24}, \textbf{25}, and \textbf{26}.
    \item Unlike the older xv6 book implementation which bundled sleeping and locking into a single call (\texttt{sleep(chan, \&lock)}), our kernel employs a decoupled two-phase sleep mechanism:
    \begin{center}
        \texttt{sleep\_prepare(chan)} $\longrightarrow$ \texttt{release(\&lock)} $\longrightarrow$ \texttt{sleep()} $\longrightarrow$ \texttt{acquire(\&lock)}
    \end{center}
    This prevents lost-wakeup race conditions while ensuring that a process never goes to sleep while holding a kernel spinlock.
\end{itemize}

% ==============================================================================
\subsection{Part 1. The numbers and the user side}

\subsubsection{Overview of User-Side Components}
Before kernel handlers are constructed, the user-space environment and build pipeline must recognize the new system call numbers and interfaces:

\begin{enumerate}[leftmargin=*]
    \item \textbf{1.1 System Call Numbers (\texttt{kernel/syscall.h}):}\\
    We assign numeric constants 24, 25, and 26 to identify each new system call:
\begin{lstlisting}[style=c]
#define SYS_sync       22
#define SYS_trace      23
#define SYS_sem_init   24
#define SYS_sem_down   25
#define SYS_sem_up     26
\end{lstlisting}

    \item \textbf{1.2 User Header Declarations (\texttt{user/user.h}):}\\
    We declare the prototypes available to user programs:
\begin{lstlisting}[style=c]
int sem_init(int, int);
int sem_down(int);
int sem_up(int);
\end{lstlisting}

    \item \textbf{1.3 Assembly Stub Generation (\texttt{user/usys.pl}):}\\
    We append entry points to \texttt{user/usys.pl}, generating machine assembly stubs that load the syscall number into register \texttt{a7} and execute \texttt{ecall}:
\begin{lstlisting}[style=bash]
entry("sem_init");
entry("sem_down");
entry("sem_up");
\end{lstlisting}

    \item \textbf{1.4 Test Application (\texttt{user/semtest.c}):}\\
    We implement \texttt{semtest.c} containing two distinct synchronization test suites:
    \begin{itemize}[noitemsep]
        \item \textbf{Test 1 (Order Enforcement):} Semaphore 0 starts at $0$. The child calls \texttt{sem\_down(0)}, forcing it to wait until the parent sleeps 10 ticks and signals \texttt{sem\_up(0)}.
        \item \textbf{Test 2 (Strict Ping-Pong Alternation):} Semaphore 1 (ping's turn) starts at $1$; Semaphore 2 (pong's turn) starts at $0$. Two processes alternate sequentially for three iterations.
    \end{itemize}

\begin{lstlisting}[style=c]
// semtest.c: two tests for the kernel semaphores (Lab 12).
#include "kernel/types.h"
#include "user/user.h"

int
main(void)
{
  // Test 1: order. Semaphore 0 starts at 0, so the child must wait.
  sem_init(0, 0);
  if (fork() == 0) {
    sem_down(0);
    printf("child: I ran after the parent\n");
    exit(0);
  }
  pause(10);
  printf("parent: done, I let the child go\n");
  sem_up(0);
  wait(0);

  // Test 2: ping-pong. Semaphore 1 = ping's turn, semaphore 2 = pong's turn.
  sem_init(1, 1);
  sem_init(2, 0);
  if (fork() == 0) {
    for (int i = 1; i <= 3; i++) {
      sem_down(2);
      printf(" pong %d\n", i);
      sem_up(1);
    }
    exit(0);
  }
  for (int i = 1; i <= 3; i++) {
    sem_down(1);
    printf("ping %d\n", i);
    sem_up(2);
  }
  wait(0);
  printf("semtest: OK\n");
  exit(0);
}
\end{lstlisting}

    \item \textbf{1.5 Multi-Core Critical Section Test (\texttt{user/mix.c}):}\\
    Two processes concurrently print five lines each, writing one character per \texttt{write()} syscall. Without synchronization, character streams interleave; with semaphore 3 enabled as a lock, output remains clean:

\begin{lstlisting}[style=c]
// mix.c: two processes print five lines each, one letter per write().
// mix no lock: the letters of the two processes can mix
// mix lock semaphore 3 used as a lock: every line stays whole
#include "kernel/types.h"
#include "user/user.h"

void
lines(char *s, int uselock)
{
  for (int k = 0; k < 5; k++) {
    if (uselock)
      sem_down(3);
    for (char *p = s; *p; p++)
      write(1, p, 1);
    if (uselock)
      sem_up(3);
  }
}

int
main(int argc, char *argv[])
{
  int uselock = (argc > 1);
  sem_init(3, 1);
  if (fork() == 0) {
    lines("bbbbbbbbbbbbbbbbbbbbbbbbbbbbbb\n", uselock);
    exit(0);
  }
  lines("AAAAAAAAAAAAAAAAAAAAAAAAAAAAAA\n", uselock);
  wait(0);
  exit(0);
}
\end{lstlisting}

    \item \textbf{1.6 Makefile Integration:}\\
    We append \texttt{\$U/\_semtest\textbackslash} and \texttt{\$U/\_mix\textbackslash} to \texttt{UPROGS}.
\end{enumerate}

\subsubsection{Check 1: Verification of Unimplemented Syscalls}
Running \texttt{semtest} prior to kernel implementation produces expected dispatch errors:
\begin{lstlisting}[style=bash]
$ semtest
3 semtest: unknown sys call 24
4 semtest: unknown sys call 25
child: I ran after the parent
parent: done, I let the child go
\end{lstlisting}
Crucially, notice that the child executed \emph{before} the parent because the kernel returned error code -1 immediately instead of suspending the child process.

% ==============================================================================
\subsection{Part 2. The kernel side}

\subsubsection{In-Kernel Semaphore Subsystem: kernel/sem.c}
We introduce a global table of 8 counting semaphores (\texttt{struct sem sems[NSEM]}). Each semaphore contains:
\begin{itemize}[noitemsep]
    \item A spinlock (\texttt{struct spinlock lock}) protecting internal updates;
    \item An integer value (\texttt{int value}) tracking available permits.
\end{itemize}

\begin{lstlisting}[style=c]
// sem.c: counting semaphores inside the xv6 kernel (Lab 12).
// A small table of NSEM semaphores, shared by every process.
// Each one is a number plus a spinlock that protects the number.
#include "types.h"
#include "riscv.h"
#include "defs.h"
#include "param.h"
#include "spinlock.h"
#include "proc.h"

#define NSEM 8

struct sem {
  struct spinlock lock; // a zeroed spinlock is a free spinlock
  int value;            // how many may still pass
};

static struct sem sems[NSEM];

// sem_init(id, value): give semaphore id its start value.
uint64
sys_sem_init(void)
{
  int id, value;
  argint(0, &id);
  argint(1, &value);
  if (id < 0 || id >= NSEM || value < 0)
    return -1;
  acquire(&sems[id].lock);
  sems[id].value = value;
  release(&sems[id].lock);
  return 0;
}

// sem_down(id): take one unit; if there is none, sleep until there is.
uint64
sys_sem_down(void)
{
  int id;
  argint(0, &id);
  if (id < 0 || id >= NSEM)
    return -1;
  struct sem *s = &sems[id];
  acquire(&s->lock);
  while (s->value == 0) {
    sleep_prepare(s);       // I will sleep on the address of s
    release(&s->lock);      // never sleep holding a spinlock
    sleep();                // returns at once if sem_up came in between
    acquire(&s->lock);
    if (killed(myproc())) { // kill must still work on a waiting process
      release(&s->lock);
      return -1;
    }
  }
  s->value--;
  release(&s->lock);
  return 0;
}

// sem_up(id): give one unit back and wake whoever sleeps on s.
uint64
sys_sem_up(void)
{
  int id;
  argint(0, &id);
  if (id < 0 || id >= NSEM)
    return -1;
  struct sem *s = &sems[id];
  acquire(&s->lock);
  s->value++;
  wakeup(s);
  release(&s->lock);
  return 0;
}
\end{lstlisting}

\subsubsection{Build System and Dispatch Table Registration}
\begin{enumerate}[label=(\alph*), noitemsep]
    \item In \texttt{Makefile}, \texttt{\$K/sem.o \textbackslash} was added to the \texttt{OBJS} kernel object list;
    \item In \texttt{kernel/syscall.c}, \texttt{sys\_sem\_init}, \texttt{sys\_sem\_down}, and \texttt{sys\_sem\_up} were registered:
\begin{lstlisting}[style=c]
extern uint64 sys_sem_init(void);
extern uint64 sys_sem_down(void);
extern uint64 sys_sem_up(void);
...
static uint64 (*syscalls[])(void) = {
...
  [SYS_sync]     = sys_sync,
  [SYS_sem_init] = sys_sem_init,
  [SYS_sem_down] = sys_sem_down,
  [SYS_sem_up]   = sys_sem_up,
};
\end{lstlisting}
\end{enumerate}

\subsubsection{Check 2: Verification of Kernel Semaphores}
We booted xv6 under QEMU and ran \texttt{semtest}:

\begin{lstlisting}[style=bash]
$ semtest
parent: done, I let the child go
child: I ran after the parent
ping 1
 pong 1
ping 2
 pong 2
ping 3
 pong 3
semtest: OK
\end{lstlisting}

The child process is strictly suspended inside \texttt{sem\_down(0)} during the parent's 10-tick sleep, verifying zero-CPU blocking and ordered wakeups. Ping and pong alternate with mathematical precision.

% ==============================================================================
\subsection{Part 3. A semaphore as a lock: mix}

\subsubsection{Interleaved Multi-Core Execution Without Locks}
Because QEMU virtualizes an SMP system with 3 hardware CPU cores (\texttt{-smp 3}), two processes writing single characters through \texttt{write(1, p, 1)} execute simultaneously in real time. Without synchronization, the character streams collide and fragment:

\begin{lstlisting}[style=bash]
$ mix
AAAAAAAbAbbAbAbbbbbAbAbAbbAbAbbbbbbbbbbbbbbbb
AbAbbbbbbbAbbAbAbAbAbAbbbbbbbbbbbbbbbAb
AAAAbbAA
bbbbbbbbbbbAbAAAAbbbAbAbAbAbAbAbAbAbAbAbAbAbAbA
AbAbAbAbAbAbAbAbAbbbAbbAb
\end{lstlisting}

\subsubsection{Enforcing Coherent Output With Kernel Locks}
Invoking \texttt{mix lock} utilizes kernel semaphore 3 (initialized to 1) to wrap the write loop in mutual exclusion:

\begin{lstlisting}[style=bash]
$ mix lock
AAAAAAAAAAAAAAAAAAAAAAAAAAAAAA
AAAAAAAAAAAAAAAAAAAAAAAAAAAAAA
AAAAAAAAAAAAAAAAAAAAAAAAAAAAAA
AAAAAAAAAAAAAAAAAAAAAAAAAAAAAA
AAAAAAAAAAAAAAAAAAAAAAAAAAAAAA
bbbbbbbbbbbbbbbbbbbbbbbbbbbbbb
bbbbbbbbbbbbbbbbbbbbbbbbbbbbbb
bbbbbbbbbbbbbbbbbbbbbbbbbbbbbb
bbbbbbbbbbbbbbbbbbbbbbbbbbbbbb
bbbbbbbbbbbbbbbbbbbbbbbbbbbbbb
\end{lstlisting}
Every line is emitted as a contiguous, uninterruptible unit.

\subsubsection*{Observe:}
\begin{enumerate}[leftmargin=*]
    \item \textbf{Without the lock, why do the letters mix? (How many system calls does one line cost?)}\\[0.2em]
    Each line consists of 30 payload letters plus 1 newline character ($\text{'\textbackslash n'}$), costing exactly \textbf{31 individual \texttt{write()} system calls per line}.\\
    Because xv6 runs across 3 separate CPU cores in parallel, the parent and child processes execute concurrently. Between any two successive single-character \texttt{write()} traps of process 1, the hardware timer interrupt or parallel execution of process 2 issues its own \texttt{write()} trap to the console UART device, interleaving the output characters on screen.

    \item \textbf{With the lock every line is whole. Which two lines of mix.c make one line a critical section?}\\[0.2em]
    Line 174:
    \begin{center}
        \texttt{if (uselock) sem\_down(3);}
    \end{center}
    and Line 178:
    \begin{center}
        \texttt{if (uselock) sem\_up(3);}
    \end{center}
    By acquiring semaphore 3 before emitting characters and releasing it only after the final newline is flushed, the entire 31-call sequence executes atomically as a single protected critical section.

    \item \textbf{In sem\_down the test is "while (s->value == 0)", not "if". Why can a process wake up and still find 0?}\\[0.2em]
    When a thread finishes in \texttt{sys\_sem\_up}, \texttt{wakeup(s)} moves \emph{all} processes sleeping on channel \texttt{s} from state \texttt{SLEEPING} to state \texttt{RUNNABLE}.\\
    However, if multiple processes were waiting, the first process scheduled acquires \texttt{s->lock}, decrements \texttt{s->value} from 1 to 0, and exits \texttt{sem\_down}. When the second awakened process subsequently acquires \texttt{s->lock}, it discovers \texttt{s->value == 0}. If the condition were evaluated using \texttt{if}, the second process would incorrectly decrement \texttt{s->value} to $-1$, violating the fundamental semaphore invariant.

    \item \textbf{sem\_down calls release(\&s->lock) BEFORE sleep(). What would go wrong if it slept holding the spinlock?}\\[0.2em]
    Holding a spinlock while sleeping causes immediate \textbf{system-wide deadlock}:
    \begin{itemize}[noitemsep]
        \item In xv6, interrupts are disabled while holding a spinlock.
        \item More critically, another process attempting to signal \texttt{sem\_up()} must acquire \texttt{\&s->lock} before incrementing \texttt{s->value} and calling \texttt{wakeup(s)}.
        \item If the waiting process went to sleep holding \texttt{s->lock}, the signaling process would spin forever in \texttt{acquire(\&s->lock)}, unable to make progress or call \texttt{wakeup()}. The sleeping process would never wake up, permanently freezing the kernel.
    \end{itemize}
\end{enumerate}

% ==============================================================================
\section{Conclusion}

Through Laboratory Work 13, I implemented an in-kernel counting semaphore facility within the **xv6-riscv** operating system.

By creating system calls 24, 25, and 26, we followed the hardware exception lifecycle from user-space assembly stubs across the RISC-V supervisor boundary into kernel dispatch routines. Investigating \texttt{sys\_sem\_down()} and \texttt{sys\_sem\_up()} illuminated the fundamental mechanics of operating system synchronization:
\begin{itemize}[noitemsep]
    \item A semaphore fundamentally pairs an integer counter with a spinlock and sleep/wakeup scheduling channels.
    \item Two-phase sleeping (\texttt{sleep\_prepare} and \texttt{sleep}) eliminates lost-wakeup race conditions without violating spinlock occupancy invariants.
    \item Multi-process testing with \texttt{semtest} and \texttt{mix} empirically demonstrated how kernel-managed semaphores eliminate busy-waiting, protect multi-core console buffers, and coordinate complex execution sequences.
\end{itemize}

This laboratory successfully completed our exploration of operating systems synchronization, uniting external POSIX library usage with internal kernel-level implementation.

\end{document}
